\documentclass[../../main.tex]{subfiles}

\begin{document}

\section{线性算子的伴随与对称性}\label{section:ch16-adjoint}

在本节中, $H$ 始终表示 Hilbert 空间. 我们记 $\mathcal{L}(H,\ H)$ 为 $\mathcal{L}(H)$. 令 $T$ 属于 $\mathcal{L}(H)$. 对于 $H$ 中的固定向量 $v$, 映射
\begin{align*}
u\mapsto\langle T(u),v\rangle,\quad u\in H
\end{align*}
属于 $H^{*}$, 由于它是线性的, 且根据 Cauchy-Schwarz 不等式, 有 $|\langle T(u),\ v\rangle|\leqslant c\cdot\|u\|$, 其中 $c=\|T\|\cdot\|v\|$. 根据 Riesz-Fr\'echet 表示定理, 存在唯一向量 $h\in H$ 使得对所有 $u\in H$, $\langle T(u),\ v\rangle=\langle h,\ u\rangle=\langle u,\ h\rangle$. 我们将此向量 $h$ 记为 $T^{*}(v)$. 这定义了映射 $T^{*}:H\to H$, 它由下面的关系式确定:
\begin{align}
\text{对所有 }u,v\in H,\quad\langle T(u),v\rangle=\langle u,T^{*}(v)\rangle.\label{align:ch16-adjoint-def}
\end{align}

我们称 $T^{*}$ 为 $T$ 的伴随.

\begin{proposition}\label{proposition:ch16-adjoint-bounded}
令 $H$ 为 Hilbert 空间. 若 $T$ 属于 $\mathcal{L}(H)$, 则 $T^{*}$ 也如此, 且 $\|T\|=\|T^{*}\|$.
\end{proposition}

\begin{proof}
显然 $T^{*}$ 是线性的. 令 $h$ 为 $H$ 中的单位向量. 则根据 Cauchy-Schwarz 不等式,
\begin{align*}
\|T^{*}(h)\|^{2}=\langle T^{*}(h),T^{*}(h)\rangle=\langle T(T^{*}(h)),h\rangle\leqslant\|T\|\|T^{*}(h)\|.
\end{align*}
因此 $T^{*}$ 属于 $\mathcal{L}(H)$ 且 $\|T^{*}\|\leqslant\|T\|$. 但也观察到
\begin{align*}
\|T(h)\|^{2}=\langle T(h),T(h)\rangle=\langle T^{*}(T(h)),h\rangle\leqslant\|T^{*}\|\|T(h)\|.
\end{align*}
因此 $\|T\|\leqslant\|T^{*}\|$.
\end{proof}

我们把证明伴随的以下结构性质留作练习. 对 $T,\ S\in\mathcal{L}(H)$,
\begin{align}
(T^{*})^{*}=T,\ (T+S)^{*}=T^{*}+S^{*}\text{ 且 }(T\circ S)^{*}=S^{*}\circ T^{*}.\label{align:ch16-adjoint-algebra}
\end{align}

\begin{proposition}\label{proposition:ch16-kernel-orth}
令 $H$ 为 Hilbert 空间. 假定 $T$ 属于 $\mathcal{L}(H)$ 且有闭的象. 则
\begin{align}
\operatorname{Im}T\oplus\ker T^{*}=H.\label{align:ch16-kernel-orth}
\end{align}
\end{proposition}

\begin{proof}
由于 $\operatorname{Im}T$ 是闭的, 根据定理\cref{theorem:ch16-orthogonal-decomposition}仅须证明 $\ker T^{*}=[\operatorname{Im}T]^{\perp}$. 而这是关系式\eqref{align:ch16-adjoint-def}的一个直接推论.
\end{proof}

我们称算子 $T\in\mathcal{L}(H)$ 是可逆的, 若它是一对一与映上的.

\begin{proposition}\label{proposition:ch16-coercive-invertible}
令 $H$ 为 Hilbert 空间. 假定 $T$ 属于 $\mathcal{L}(H)$ 且存在 $c>0$ 使得
\begin{align}
\text{对所有 }h\in H,\quad\langle T(h),h\rangle\geqslant c\cdot\|h\|^{2}.\label{align:ch16-coercive}
\end{align}
则 $T$ 是可逆的.
\end{proposition}

\begin{proof}
不等式\eqref{align:ch16-coercive}蕴涵 $\ker T=\{0\}$. 我们从\eqref{align:ch16-coercive}和 Cauchy-Schwarz 不等式推出
\begin{align*}
\text{对所有 }u,v\in H,\quad\|T(u)-T(v)\|\geqslant c\|u-v\|.
\end{align*}
我们宣称 $T$ 有闭的象. 事实上, 令在 $H$ 中强意义下 $\{T(h_{n})\}\to h$. 则 $\{T(h_{n})\}$ 是 Cauchy 的. 以上不等式蕴涵 $\{h_{n}\}$ 是 Cauchy 的. 但 $H$ 是完备的, 因此存在向量 $h_{*}$ 使得 $\{h_{n}\}$ 强收敛于它. 根据 $T$ 的连续性, $T(h_{*})=h$. 因此 $T(H)$ 是闭的. 我们也宣称 $\ker T^{*}=\{0\}$. 事实上, 根据\eqref{align:ch16-coercive}、内积的对称性以及伴随的定义,
\begin{align*}
\text{对所有 }h\in H,\quad\langle T^{*}(h),h\rangle=\langle T(h),h\rangle\geqslant c\cdot\|h\|^{2}.
\end{align*}
因此 $\ker T^{*}=\{0\}$. 我们从前一个命题推出 $T(H)=H$.
\end{proof}

检查 Riesz-Fr\'echet 表示定理的证明发现没有用到内积的对称性. 该定理的以下重要推广在偏微分方程的研究中有很多应用.

\begin{theorem}[Lax-Milgram 引理]\label{theorem:ch16-lax-milgram}
令 $H$ 为 Hilbert 空间. 假定函数 $B:H\times H\to\mathbf{R}$ 有下面三个性质:
\begin{enumerate}[(i)]
\item 对每个 $u\in H$, $H$ 上的以下两个泛函是线性的: $v\mapsto B(u,\ v)$ 与 $v\mapsto B(v,\ u)$.
\item 存在 $c_{1}>0$ 使得
\begin{align*}
\text{对所有 }u,v\in H,\quad|B(u,v)|\leqslant c_{1}\cdot\|u\|\cdot\|v\|.
\end{align*}
\item 存在 $c_{2}>0$ 使得
\begin{align*}
\text{对所有 }h\in H,\quad B(h,h)\geqslant c_{2}\cdot\|h\|^{2}.
\end{align*}
\end{enumerate}
则对每个 $\psi\in H^{*}$, 存在唯一的 $h\in H$ 使得
\begin{align*}
\text{对所有 }u\in H,\quad\psi(u)=B(h,u).
\end{align*}
\end{theorem}

\begin{proof}
令 $T:H\to H^{*}$ 为 Riesz-Fr\'echet 表示定理定义的同构, 即对每个 $h\in H$,
\begin{align}
T(h)[u]=\langle h,u\rangle.\label{align:ch16-lax-1}
\end{align}
对每个 $h$, 定义泛函 $S(h):H\to\mathbf{R}$ 为
\begin{align}
\text{对所有 }u\in H,\quad S(h)[u]=B(h,u).\label{align:ch16-lax-2}
\end{align}
我们从假设(i)和(ii)推出每个 $S(h)$ 是 $H$ 上的有界线性泛函. 且算子 $S:H\to H^{*}$ 是线性与连续的. 由于 $T$ 是 $H$ 映上 $H^{*}$ 的同构, 则证明 $S$ 是 $H$ 映上 $H^{*}$ 的同构等价于证明算子 $T^{-1}\circ S\in\mathcal{L}(H)$ 是可逆的. 然而, 根据假设(iii),
\begin{align*}
\text{对所有 }h\in H,\quad\langle(T^{-1}\circ S)(h),h\rangle=S(h)[h]=B(h,h)\geqslant c_{2}\cdot\|h\|^{2}.
\end{align*}
前一个命题告诉我们 $T^{-1}\circ S\in\mathcal{L}(H)$ 是可逆的.
\end{proof}

\begin{definition}[对称/自伴算子]\label{definition:ch16-symmetric}
算子 $T\in\mathcal{L}(H)$ 称为是对称的或自伴的, 若 $T=T^{*}$, 即
\begin{align*}
\text{对所有 }u,v\in H,\quad\langle T(u),v\rangle=\langle u,T(v)\rangle.
\end{align*}
\end{definition}

\begin{example}\label{example:ch16-symmetric-matrix}
令 $\{\varphi_{k}\}$ 为可分 Hilbert 空间 $H$ 的规范正交基而 $T$ 属于 $\mathcal{L}(H)$. 则根据内积的连续性, $T$ 是对称的当且仅当
\begin{align*}
\text{对所有 }1\leqslant i,j<\infty,\quad\langle T(\varphi_{i}),\varphi_{j}\rangle=\langle T(\varphi_{j}),\varphi_{i}\rangle.
\end{align*}
特别地, 若 $H$ 是欧氏空间 $\mathbf{R}^{n}$, 则 $T$ 是对称的当且仅当在规范正交基下表示 $T$ 的 $n\times n$ 矩阵是对称矩阵.
\end{example}

对称算子 $T\in\mathcal{L}(H)$ 称为是非负的, 写为 $T\geqslant 0$, 若对所有 $h\in H$, $\langle T(h),\ h\rangle\geqslant 0$. 此外, 对两个对称算子 $A,\ B\in\mathcal{L}(H)$, 我们写 $A\geqslant B$ 意味着 $A-B\geqslant 0$. 非负对称算子的和是非负的与对称的. 此外,
\begin{align}
\text{若 }T\in\mathcal{L}(H)\text{ 是对称与非负的, 则对任何 }S\in\mathcal{L}(H),S^{*}TS\text{ 也是如此}.\label{align:ch16-nonneg-star}
\end{align}
这是由于对每个 $h\in H$, $\langle S^{*}TS(h),\ h\rangle=\langle TS(h),\ S(h)\rangle\geqslant 0$. 在习题\cref{exercise:ch16-37}--\cref{exercise:ch16-43}中, 我们将探讨对称算子中的这个序关系的一些有趣推论.

\begin{theorem}[极化等式]\label{theorem:ch16-polarization}
对于对称算子 $T\in\mathcal{L}(H)$,
\begin{align}
\text{对所有 }u,v\in H,\quad\langle T(u),v\rangle=\frac{1}{4}[\langle T(u+v),u+v\rangle-\langle T(u-v),u-v\rangle].\label{align:ch16-polarization}
\end{align}
\end{theorem}

为证明这个等式, 简单地将右边的两个内积展开即可. 若我们把算子 $T\in\mathcal{L}(H)$ 与对所有 $u\in H$ 定义为 $Q_{T}(u)=\langle T(u),\ u\rangle$ 的二次型 $Q_{T}:H\to\mathbf{R}$ 相联系, 极化等式告诉我们 $T$ 由 $Q_{T}$ 完全确定. 特别地, 在 $H$ 上 $T=0$ 当且仅当在 $H$ 上 $Q_{T}=0$. 事实上, 以下更精确的结果成立. 对 $T\in\mathcal{L}(H)$ 与 $\lambda\in\mathbf{R}$, 记 $\lambda\operatorname{Id}-T$ 为 $\lambda-T$, 其中恒等映射 $\operatorname{Id}:H\to H$ 定义为对所有 $h$, $\operatorname{Id}(h)=h$.

\begin{proposition}\label{proposition:ch16-norm-form}
令 $H$ 为 Hilbert 空间而算子 $T\in\mathcal{L}(H)$ 是对称的. 则
\begin{align}
\|T\|=\sup_{\|u\|=1}|\langle Tu,u\rangle|.\label{align:ch16-norm-form}
\end{align}
\end{proposition}

\begin{proof}
记 $\sup_{\|u\|=1}|\langle T(u),\ u\rangle|$ 为 $\eta$. 若 $\eta=0$, 我们从极化等式推出 $T=0$. 因此考虑 $\eta>0$ 的情形. 观察到, 根据 Cauchy Schwarz 不等式, 对单位向量 $u\in H$,
\begin{align*}
|\langle T(u),u\rangle|\leqslant\|T(u)\|\|u\|\leqslant\|T\|.
\end{align*}
因此 $\eta\leqslant\|T\|$. 为证明反向的不等式, 观察到对称算子 $\eta-T$ 与 $\eta+T$ 是非负的, 因此根据\eqref{align:ch16-nonneg-star}, 算子
\begin{align*}
(\eta+T)^{*}(\eta-T)(\eta+T)=(\eta+T)(\eta-T)(\eta+T)
\end{align*}
与
\begin{align*}
(\eta-T)^{*}(\eta+T)(\eta-T)=(\eta-T)(\eta+T)(\eta-T)
\end{align*}
也是非负的, 因此它们的和 $2\eta(\eta^{2}-T^{2})$ 也如此. 由于 $2\eta>0$, $\eta^{2}-T^{2}$ 是非负的, 即
\begin{align*}
\text{对所有 }u\in H,\quad\|T(u)\|^{2}=\langle T(u),T(u)\rangle=\langle T^{2}(u),u\rangle\leqslant\eta^{2}\langle u,u\rangle=\eta^{2}\|u\|^{2}.
\end{align*}
因此 $\|T\|\leqslant\eta$.
\end{proof}

在线性算子 $T\in\mathcal{L}(H)$ 的研究中, 一个一般的策略是将 $H$ 表示为直和 $H_{1}\oplus H_{2}$, 使得 $T(H_{1})\subseteq H_{1}$ 且 $T(H_{2})\subseteq H_{2}$. 在这种情况下我们说分解 $H=H_{1}\oplus H_{2}$ 约化算子 $T$. 一般地, 若 $T(H_{1})\subseteq H_{1}$, 我们不能推出 $T(H_{2})\subseteq H_{2}$. 然而, 对 $H$ 上的对称算子与 $H$ 的正交直和分解, 我们有以下简单但非常有用的结果.

\begin{proposition}\label{proposition:ch16-invariant-orthocomp}
令 $H$ 为 Hilbert 空间. 假定算子 $T\in\mathcal{L}(H)$ 是对称的而 $V$ 是 $H$ 的子空间使得 $T(V)\subseteq V$. 则 $T(V^{\perp})\subseteq V^{\perp}$.
\end{proposition}

\begin{proof}
令 $u$ 属于 $V^{\perp}$. 则对任何 $v\in V$, $\langle T(u),\ v\rangle=\langle u,\ T(v)\rangle$, 而 $\langle u,\ T(v)\rangle=0$, 由于 $T(V)\subseteq V$ 且 $u\in V^{\perp}$. 因此 $T(u)\in V^{\perp}$.
\end{proof}

\begin{exercise}\label{exercise:ch16-27}
证明\eqref{align:ch16-adjoint-algebra}.
\end{exercise}

\begin{exercise}\label{exercise:ch16-28}
令 $T$ 和 $S$ 属于 $\mathcal{L}(H)$ 且是对称的. 证明 $T=S$ 当且仅当 $Q_{T}=Q_{S}$.
\end{exercise}

\begin{exercise}\label{exercise:ch16-29}
证明对称算子全体是 $\mathcal{L}(H)$ 的子空间. 也证明, 若 $S$ 和 $T$ 是对称的, 则它们的复合 $S\circ T$ 是对称的当且仅当 $T$ 与 $S$ 可交换, 即 $S\circ T=T\circ S$.
\end{exercise}

\begin{exercise}\label{exercise:ch16-30}
(Hellinger-Toeplitz)令 $H$ 为 Hilbert 空间而线性算子 $T:H\to H$ 具有性质: 对所有 $u,\ v\in H$, $\langle T(u),\ v\rangle=\langle u,\ T(v)\rangle$. 证明 $T$ 属于 $\mathcal{L}(H)$.
\end{exercise}

\begin{exercise}\label{exercise:ch16-31}
展示一个算子 $T\in\mathcal{L}(\mathbf{R}^{2})$ 使得 $\|T\|>\sup_{\|u\|=1}|\langle T(u),\ u\rangle|$.
\end{exercise}

\begin{exercise}\label{exercise:ch16-32}
令 $\mathcal{L}(H)$ 中的 $T$ 和 $S$ 是对称的. 假设 $S\geqslant T$ 且 $T\geqslant S$. 证明 $T=S$.
\end{exercise}

\begin{exercise}\label{exercise:ch16-33}
令 $V$ 为 Hilbert 空间的非平凡闭子空间, 而 $P$ 是 $H$ 映上 $V$ 的正交投影. 证明 $P=P^{*},\ P\geqslant 0$, 且 $\|P\|=1$.
\end{exercise}

\begin{exercise}\label{exercise:ch16-34}
令 $P\in\mathcal{L}(H)$ 是投影. 证明 $P$ 是 $H$ 映上 $P(H)$ 的正交投影当且仅当 $P=P^{*}$.
\end{exercise}

\begin{exercise}\label{exercise:ch16-35}
令 $\{\varphi_{k}\}$ 为 Hilbert 空间 $H$ 的规范正交基且对每个自然数 $n$, 定义 $P_{n}$ 为 $H$ 映上 $\{\varphi_{1},\ \cdots,\ \varphi_{n}\}$ 的线性扩张的正交投影. 证明 $P_{n}$ 是对称的且
\begin{align*}
\text{对所有 }n,\quad 0\leqslant P_{n}\leqslant P_{n+1}\leqslant\operatorname{Id}.
\end{align*}
证明在 $H$ 上 $\{P_{n}\}$ 逐点收敛到 $\operatorname{Id}$, 但在单位球上不一致收敛.
\end{exercise}

\begin{exercise}\label{exercise:ch16-36}
证明: 若 $T\in\mathcal{L}(H)$ 是可逆的, $T^{*}\circ T$ 也如此从而 $T^{*}$ 也如此.
\end{exercise}

\begin{exercise}\label{exercise:ch16-37}
(一般的 Cauchy-Schwarz 不等式)令 $T\in\mathcal{L}(H)$ 为对称与非负的. 证明对所有 $u,\ v\in H$,
\begin{align*}
|\langle T(u),v\rangle|^{2}\leqslant\langle T(u),u\rangle\cdot\langle T(v),v\rangle.
\end{align*}
\end{exercise}

\begin{exercise}\label{exercise:ch16-38}
用前一个习题证明: 若 $S,\ T\in\mathcal{L}(H)$ 是对称的且 $S\geqslant T$, 则对每个 $u\in H$,
\begin{align*}
\|S(u)-T(u)\|^{4}&=\langle(S-T)(u),(S-T)(u)\rangle^{2}\leqslant|\langle(S-T)(u),u\rangle|\langle(S-T)^{2}(u),(S-T)(u)\rangle|.
\end{align*}
因此得出
\begin{align*}
\|S(u)-T(u)\|^{4}\leqslant|\langle S(u),u\rangle-\langle T(u),u\rangle|\cdot\|S-T\|^{3}\cdot\|u\|^{2}.
\end{align*}
\end{exercise}

\begin{exercise}\label{exercise:ch16-39}
(关于对称算子的单调收敛定理)$\mathcal{L}(H)$ 中的对称算子序列 $\{T_{n}\}$ 称为是单调递增的, 若对每个 $n$, $T_{n+1}\geqslant T_{n}$; 称为是有上界的, 若存在 $\mathcal{L}(H)$ 中的对称算子 $S$ 使得对所有 $n$, $T_{n}\leqslant S$.
\begin{enumerate}[(i)]
\item 用前一个习题证明 $\mathcal{L}(H)$ 中的单调递增的对称算子序列 $\{T_{n}\}$ 逐点收敛于 $\mathcal{L}(H)$ 中的对称算子当且仅当它有上界.
\item 证明 $\mathcal{L}(H)$ 中的对称算子的单调递增序列 $\{T_{n}\}$ 有上界当且仅当它是逐点有界的, 即对每个 $h\in H$, 序列 $\{T_{n}(h)\}$ 是有界的.
\end{enumerate}
\end{exercise}

\begin{exercise}\label{exercise:ch16-40}
令 $S\in\mathcal{L}(H)$ 为非负对称算子使得 $0\leqslant S\leqslant\operatorname{Id}$. 令 $T_{1}=1/2(\operatorname{Id}-S)$, 而若 $n$ 是自然数使得 $T_{n}\in\mathcal{L}(H)$ 被定义, 定义 $T_{n+1}=1/2(\operatorname{Id}-S+T_{n}^{2})$, 如此定义 $\mathcal{L}(H)$ 中的算子序列 $\{T_{n}\}$.
\begin{enumerate}[(i)]
\item 证明对每个自然数 $n$, $T_{n}$ 与 $T_{n+1}-T_{n}$ 是具有非负系数的 $\operatorname{Id}-S$ 的多项式.
\item 证明 $\{T_{n}\}$ 是对称算子的单调递增序列, 它的上界为 $\operatorname{Id}$.
\item 用前一个习题证明 $\{T_{n}\}$ 逐点收敛于对称算子 $T$, 使得 $0\leqslant T\leqslant\operatorname{Id}$ 且 $T=1/2(\operatorname{Id}-S+T^{2})$.
\item 定义 $A=(\operatorname{Id}-T)$. 证明 $A^{2}=S$.
\end{enumerate}
\end{exercise}

\begin{exercise}\label{exercise:ch16-41}
(非负对称算子的平方根)令 $T\in\mathcal{L}(H)$ 为非负对称算子. 非负对称算子 $A\in\mathcal{L}(H)$ 称为是 $T$ 的平方根, 若 $A^{2}=T$. 用前一个习题的归纳构造证明 $T$ 有平方根 $A$, 它与 $\mathcal{L}(H)$ 中每个与 $T$ 可交换的算子可交换. 证明平方根是唯一的: 记为 $\sqrt{T}$. 最后, 证明 $T$ 是可逆的当且仅当 $\sqrt{T}$ 是可逆的.
\end{exercise}

\begin{exercise}\label{exercise:ch16-42}
可逆算子 $T\in\mathcal{L}(H)$ 称为是正交的, 若 $T^{-1}=T^{*}$. 证明可逆算子是正交的当且仅当它是一个等距.
\end{exercise}

\begin{exercise}\label{exercise:ch16-43}
(极分解)令 $T\in\mathcal{L}(H)$ 是可逆的. 证明存在正交可逆算子 $A$ 与非负对称可逆算子 $B$ 使得 $T=B\circ A$. (提示: 证明 $T^{*}T$ 是可逆与对称的, 且令 $B=\sqrt{T\circ T^{*}}$.)
\end{exercise}

\end{document}